\section{Compact Appendix B: interface forcing, unique interior, and conservation}
\label{app:interface}

This appendix records the minimum interface theorems the body uses once the kernel floor of Appendix~\ref{app:compactA} is fixed. Its role is not to reopen the foundational argument but to state, in compressed form, the forced interface consequences of that kernel.

\begin{definition}[Fixed-scope interface]
At fixed scope $D$, admissibility is a bivalent gate on the evaluated fragment. Standing is the positive side of that gate as it survives licensed same-scope reuse. Any same-scope parameter that would further refine admission, rejection, or standing-bearing continuation is counted as hidden gate work.
\end{definition}

\begin{theorem}[AMetric and bivalent interface forcing]
At the boundary there is no admissibility-relevant metric, selector, ranking, grading, or deferred-entry parameter. On the evaluated side, every admissibility-relevant status collapses extensionally to admitted or not admitted.
\end{theorem}

\begin{proof}
If an admissibility-relevant selector or grading existed at the boundary, admissibility would depend on pre-boundary parameterization rather than on the gate itself. If an evaluated-side distinction did more than split admitted from not admitted, it would either create a second same-scope gate or be inert bookkeeping. The first possibility is excluded by kernel uniqueness and the second carries no governance force. Hence the interface is AMetric at the boundary and bivalent on the evaluated side.
\end{proof}

\begin{theorem}[Standing emergence]
At fixed scope, standing is exactly reuse-stable admissibility. There is no independent positive-side classifier beyond the gate together with the requirement that licensed same-scope continuation preserve the same standing-bearing content.
\end{theorem}

\begin{proof}
If standing were stricter than reuse-stable admissibility, admitted acts would be split into a merely admitted class and a genuinely standing-positive class. If that distinction changed what could later be reused or continued, it would introduce hidden positive-side gate work; if it changed nothing, it would be inert bookkeeping. Therefore standing is just the positive side of the gate as preserved through licensed same-scope continuation.
\end{proof}

\begin{corollary}[No repair, no generator, no carrier]
Within the fixed scope, there is no same-act repair, no same-domain generator of standing from non-standing, and no independent carrier of inferential authority beside the gate.
\end{corollary}

\begin{proof}
A repair would reclassify the same evaluated act after failure; a generator would produce standing from a failed act without fresh gating; and a carrier would act as a second authorizer beside the gate. Each move either reopens the hidden-gate problem or violates the fixed-scope irrecoverability discipline.
\end{proof}

\begin{theorem}[Unique admissible interior and standing conservation]
At fixed scope there is exactly one admissible interior up to lawful redescription. Standing can be preserved or destroyed, but not created, amplified, transferred, or recovered inside the regime.
\end{theorem}

\begin{proof}
Plural same-scope interiors would differ by an admissibility-relevant distinction and so would require either a second gate or a parameterized boundary. Those possibilities are excluded by the preceding theorems. Conservation follows because every alleged creation, amplification, transfer, or recovery of standing would have to proceed either by same-act repair, by a generator, or by a carrier; all three have already been excluded.
\end{proof}
